\source{Лекции, с. 7–9, 57; элементарные детали доказательств восстановлены составителем}
\begin{definition}[Вершинная двусвязность рёбер]
Два различных ребра \(ab,cd\) вершинно двусвязны, если их концы можно
попарно соединить двумя путями без общих вершин, не используя сами
эти рёбра. Разрешены оба сопоставления концов; для смежных рёбер один
путь может быть нулевым. Каждое ребро считаем двусвязным с самим собой.
Эквивалентно: два различных ребра лежат на общем простом цикле.
Действительно, удаление двух рёбер из цикла даёт нужные пути, а их
склейка с этими рёбрами даёт цикл.
\end{definition}
\textbf{Элементарные детали, используемые далее.}
Граф на \(n\ge3\) вершинах назовём двусвязным, если он связен
и удаление любой одной вершины оставляет его связным.
Следующая лемма раскрывает рисунки о непересекающихся путях без обращения к потокам.
\begin{lemma}[Два выхода к циклу]\label{lem:vertex-paths}
В двусвязном графе из вершины \(x\) к множеству \(S\), \(|S|\ge2\),
\(x\notin S\), идут два пути с разными концами в \(S\), пересекающиеся
только в \(x\) и не встречающие \(S\) внутренними вершинами.
Любые две вершины двусвязного графа лежат на общем простом цикле.
\end{lemma}
\begin{proof}
В связном графе сначала построим дерево из пути от \(x\) к \(S\)
и пути от другой вершины \(S\) к первому пути. После обрезания лишнего
это три ветви от вершины \(z\): к \(x\), \(s_1\), \(s_2\), причём
\(s_1,s_2\in S\), \(s_1\ne s_2\). Одна ветвь может быть нулевой.
Ветви обрезаем при первой встрече с \(S\), так что внутри ветви \(xz\)
нет вершин \(S\). Среди таких деревьев минимизируем длину ветви \(xz\).
Если \(z\ne x\), в \(G-z\) есть путь от \(x\) к одной из двух остальных
ветвей, так как хотя бы одна из \(s_1,s_2\) отлична от \(z\).
Берём его участок от последней встречи \(t\) с \(xz\) до первой
встречи с другой ветвью. Новый обход, ветвящийся в \(t\), даёт две
разные конечные вершины \(S\) и более короткую общую ветвь \(xt\).
Противоречие. Значит, \(z=x\); получены требуемые пути.

Для второго утверждения сначала построим цикл через \(u\): у \(u\)
есть два разных соседа, соединённых простым путём в \(G-u\).
Если \(v\) вне этого цикла, применим первую часть с \(x=v\), \(S=V(C)\).
Два выхода к \(a,b\in C\) вместе с той дугой \(C\), которая содержит
\(u\), образуют простой цикл через \(u,v\).
\end{proof}
\begin{lemma}[Циклы через рёбра]\label{lem:blocks}
В двусвязном графе любые два различных ребра лежат на общем простом цикле.
Объединение двух двусвязных подграфов с двумя общими вершинами двусвязно.
\end{lemma}
\begin{proof}
Последнее утверждение следует из того, что после удаления одной вершины
оба подграфа связны и у них остаётся общая вершина.
Для первого подразделим два выбранных ребра новыми вершинами \(x,y\).
Подразделение сохраняет двусвязность: удаление старой вершины оставляет
новую прикреплённой к другому концу ребра, а удаление новой вершины
равносильно удалению немостового ребра. В двусвязном графе мостов нет:
при \(n\ge3\) хотя бы один конец моста был бы разделяющей вершиной.
По предыдущей лемме \(x,y\) лежат на общем цикле. Убирая подразделения,
получаем цикл через оба исходных ребра.
\end{proof}
\begin{theorem}[Отношение и его классы]
Вершинная двусвязность является отношением эквивалентности на \emph{рёбрах}.
Её классы называются компонентами вершинной двусвязности, или блоками;
сам блок рассматривается вместе со всеми концами его рёбер.
\end{theorem}
\begin{proof}
Рефлексивность задана, симметричность очевидна. Пусть различные пары
\(e,f\) и \(f,g\) лежат на циклах \(C,D\). Циклы двусвязны, у них общими
являются оба конца \(f\), поэтому \(C\cup D\) двусвязен.
По предыдущей лемме \(e,g\) лежат на общем цикле. Если какие-то рёбра
совпали, транзитивность очевидна. Следовательно, классы разбивают \(E\).
Немостовые классы — максимальные двусвязные подграфы: объединяем
циклы через фиксированное ребро класса, применяя лемму об объединении;
обратно, все рёбра двусвязного подграфа принадлежат одному классу.
Мост образует одноэлементный класс.
\end{proof}
Для удобства разложения графа будем также называть изолированную вершину
вырожденным блоком. Это дополнительное соглашение: класса рёбер у неё нет.
Блоки-мосты также называем вырожденными; остальные блоки двусвязны
и называются невырожденными.
Разные блоки имеют не более одной общей вершины, иначе их объединение
двусвязно и классы совпадают. Каждое ребро принадлежит ровно одному блоку.
\begin{definition}[Точка сочленения]
Вершина, принадлежащая нескольким блокам с рёбрами, называется точкой сочленения.
\end{definition}
\begin{theorem}[Критерий точки сочленения]
Вершина \(v\) является точкой сочленения тогда и только тогда,
когда её удаление увеличивает число компонент связности.
\end{theorem}
\begin{proof}
Достаточно рассмотреть компоненту, содержащую \(v\).
Если два инцидентных \(v\) ребра \(va,vb\) принадлежат разным блокам,
в \(G-v\) нет пути от \(a\) к \(b\): простой такой путь вместе
с \(av,vb\) дал бы общий цикл. Поэтому удаление разъединяет компоненту.
Обратно, если она распалась, у \(v\) есть соседи \(a,b\) в разных
полученных компонентах. Общий цикл через \(va,vb\) дал бы обходной
путь от \(a\) к \(b\) без \(v\), что невозможно. Значит, блоки разные.
Изолированная вершина и конец висячего ребра проверяются непосредственно.
\end{proof}
\begin{cor}[Критерий двусвязности]
При \(n\ge3\) эквивалентны: связность без точек сочленения;
связность после удаления любой вершины; наличие общего простого цикла
для любых двух различных вершин.
\end{cor}
\begin{proof}
Первую эквивалентность даёт критерий удаления. Вторая следует из леммы
о двух выходах; в обратную сторону цикл сохраняет связь двух оставшихся
вершин после удаления любой третьей.
\end{proof}
На вершинах отношение «лежать в одном блоке» не транзитивно:
два треугольника с общей вершиной дают контрпример. Это согласуется
с лекционным определением отношения именно на рёбрах.
\begin{lemma}[Дерево блоков]\label{lem:block-tree}
Для связного графа двудольный граф инцидентности блоков и точек
сочленения является деревом. В частности, если блоков больше одного,
есть не менее двух крайних блоков, каждый из которых содержит ровно
одну точку сочленения исходного графа.
\end{lemma}
\begin{proof}
Сначала построим граф инцидентности \(T\) всех блоков и \emph{всех}
вершин исходного графа. Ребро в \(T\) означает принадлежность вершины блоку.
Он связен: любое ребро исходного графа лежит в блоке, поэтому путь
исходного графа превращается в чередующийся маршрут по \(T\).

Пусть в \(T\) есть простой цикл. Рассмотрим объединение его блоков;
в каждом отмечены две различные вершины стыка с соседними блоками.
После удаления любой одной вершины \(x\) каждый невырожденный блок
остаётся связным; блок-мост либо сохраняется, либо от него остаётся
одна вершина стыка. Блок из изолированной вершины в цикл входить не может.
Среди различных отмеченных стыков исчезает не более одного,
поэтому вдоль остальных стыков блоки всё ещё соединяются в цепочку
или цикл. Следовательно, всё объединение без \(x\) связно.
Это двусвязный подграф, содержащий несколько различных блоков,
вопреки их максимальности (а мост вообще не может лежать в двусвязном подграфе).
Значит, \(T\) — дерево.

Если вершина \(v\) принадлежит двум или более блокам, её узел в \(T\)
имеет степень не менее двух. Удаление этого узла разъединяет дерево.
Путь в \(G-v\) между вершинами разных ветвей дал бы маршрут в
\(T-v\), что невозможно. Значит, \(v\) — точка сочленения.
Если же \(v\) принадлежит единственному блоку, его удаление не нарушает
связности остального графа: невырожденный блок остаётся связным,
а в блоке-мосте такая вершина — лист исходного графа.
Одноточечный граф рассматривается непосредственно.

Удалим из \(T\) все листовые узлы обычных вершин; останется именно
граф инцидентности блоков и точек сочленения. Он является деревом
(при единственном блоке — одной вершиной).
Дерево с более чем одной вершиной имеет не менее двух листьев:
это концы максимального простого пути. Узел точки сочленения
имеет степень не менее двух, поэтому листья соответствуют блокам.
\end{proof}
